% Full independent written audits: root and dieudonne; independent exact-integer certificate rerun retained.
\begin{lemma}[Proved: finite-length linear cyclic Bocksteins]
\label{C:ramified_linear_Bockstein}
Let $A$ be a characteristic-zero DVR of residue characteristic
$p$, let $J$ be flat over $A$, and let $k=A\ltimes J$.
Every $A/p$-linear map $\ell:J/pJ\to A/p$ induces an odd
endomorphism $b_\ell$ of the two-periodic Tate cohomology of the
reduction to $A$ of each finite free $k[C_p]$-module.
For every monic irreducible $q\in(A/\pi)[z]$, $q\ne z$, the
composition length of the $(\pi,q)$-primary even part for
$z=b_\ell^2$, taken modulo two, defines an additive invariant on
$R_k(C_p)$. It vanishes on constant representations from $A$.
\end{lemma}
\begin{proof}
Write the generator as $T_0+U$ over the split square-zero ring.
Linearize the two Tate differentials $T-1$ and
$1+T+\cdots+T^{p-1}$. Their linear terms form an odd chain map
from the Tate complex over $A$ to its tensor product with $J$.
Flatness gives the natural cohomology comparison
\[
 \widehat H^i(C_p,M\otimes_AJ)
 =\widehat H^i(C_p,M)\otimes_AJ
 =\widehat H^i(C_p,M)\otimes_{A/p}(J/pJ),
\]
because Tate cohomology is killed by $p$. Composing with $\ell$
defines $b_\ell$.
The same linear differential identities as in
Lemma~\ref{C:cyclic_linear_Bockstein} show that changes of lift or
basis induce homotopies, morphisms commute with $b_\ell$, and
exact-sequence connecting maps anticommute with $b_\ell$.
These identities use the square-zero decomposition and flatness,
not that $p$ generates the maximal ideal.

All Tate groups here have finite length over $A$. Regard them as
finite-length $A[z]$-modules via $z=b_\ell^2$. The exact
six-cycle decomposes into its maximal-ideal primary parts, and
on a part with maximal ideal $(\pi,q)$, $q\ne z$, the operator
$b_\ell$ is invertible. Folding the even and odd terms yields an
exact three-cycle; the wraparound differs by the unit $-z$, so
kernels and images are unchanged. The sum of the three composition
lengths is even. Hence the even primary length is additive
modulo two. This is length over the local ring
$A[z]_{(\pi,q)}$; a simple residue polynomial of degree $d$ has
length one, not $d$. A constant representation has zero linear
term and therefore no nonzero primary part.
\end{proof}

\begin{theorem}[Proved: flat square-zero directions at minimal cyclotomic ramification]
\label{C:cyclotomic_flat_criterion}
Let $p$ be odd, let $A$ be a characteristic-zero DVR containing a
primitive $p$-th root $\zeta$ for which $\pi=\zeta-1$ is a
uniformizer, and let $J$ be any flat $A$-module. Then
\[
 \begin{aligned}
 \vartheta_{A\ltimes J}\text{ is an isomorphism}
 &\quad\Longleftrightarrow\quad
 \vartheta_{A\ltimes J}\text{ is surjective in degree }p\\
 &\quad\Longleftrightarrow\quad pJ=J.
 \end{aligned}
\]
If $pJ\ne J$, an explicit rank-$p!$ representation of
$\mathfrak S_p$ detects failure. No finite-generation assumption
on $J$ is required.
\end{theorem}
\begin{proof}
Here $p$ is a unit times $\pi^{p-1}$. If $pJ\ne J$, also
$\pi J\ne J$; choose $x\notin\pi J$.
The module $J/pJ$ is flat over the artinian local ring $A/p$,
and is free even when its rank is infinite. For completeness,
lift a residue-field basis to obtain a map from a free module;
its cokernel equals its maximal-ideal multiple and is zero by
iteration of the nilpotent maximal ideal. Flatness then shows
that its kernel also equals its maximal-ideal multiple, so the
map is an isomorphism. Some coordinate of $x$ is a unit, and a
coordinate functional rescaled by that unit gives
$\ell:J/pJ\to A/p$ with $\ell(x)=1$.

Put
\[
 Q(t)=\frac{t^p-1}{t-\zeta},\qquad
 R(t)=\frac{Q(t)}{t-1},\qquad L=A[t]/Q(t).
\]
The polynomials $Q,R$ are monic of degrees $p-1,p-2$.
On $L$, the kernel of $t-1$ is $AR$, while
\[
 (1+t+\cdots+t^{p-1})f=(1-\zeta)f(1)R.
\]
The kernel of the latter map consists exactly of the elements
represented by polynomials divisible by $t-1$. Thus
\[
 \widehat H^0(C_p,L)=A/\pi,\qquad
 \widehat H^1(C_p,L)=0.
\]
The character line $Z$ with action $\zeta$ has the opposite
Tate groups: $H^0=0$ and $H^1=A/\pi$.

Let $a:A\to L$ send $1$ to $1$, and let $b:L\to A$ take the
coefficient of $t^{p-2}$. On $L\oplus Z$ define
\[
 T_0=\begin{pmatrix}t&0\\0&\zeta\end{pmatrix},\qquad
 U=\begin{pmatrix}0&a\\b&0\end{pmatrix},\qquad T=T_0+xU.
\]
Both off-diagonal derivatives of $T^p$ are zero: the relevant
polynomial is $Q(t)$, which annihilates $L$. Hence $T^p=1$.
The first linear Tate map sends the even generator $R$ to
$b(R)=1$ in $Z/\pi Z$. The norm derivative in the other direction
is multiplication by
\[
 \frac{(1+t+\cdots+t^{p-1})
             -(1+\zeta+\cdots+\zeta^{p-1})}{t-\zeta}
 =R(t),
\]
so it sends the odd generator $1$ to $R$. Since $\ell(x)=1$,
both Bocksteins are one. The $(z-1)$-primary even length is
therefore one.

For a generator twist $g\mapsto g^a$, $1\le a<p$, the first
linear map is multiplied by $a$ on Tate cohomology and the norm
map is unchanged. Thus $b_\ell^2=a$, and exactly $a=1$
contributes to this primary invariant. Mackey restriction after
induction to $\mathfrak S_p$ has each normalizer twist once and
otherwise only regular actions, whose Tate groups vanish.
Every Young permutation class is constant over $A$, so the
induced class is outside their span. Since $A\ltimes J$ is local,
$K_0(A\ltimes J)=\mathbb Z$ and that span is the polynomial image.
The induced rank is $p(p-1)!=p!$.

If $pJ=J$, flatness makes multiplication by $p$ bijective.
Every other prime is a unit in $A$, so
Proposition~\ref{C:divisible_nil_invariance} and
Proposition~\ref{C:DVR} give the full polynomial answer.
\end{proof}
